\documentclass[11pt,a4paper]{article}
\usepackage[utf8]{inputenc}
\usepackage[T1]{fontenc}
\usepackage[spanish,es-minimal,es-noshorthands]{babel}
\AtBeginDocument{\renewcommand{\sin}{\operatorname{sen}}\renewcommand{\arcsin}{\operatorname{arcsen}}}
\usepackage{lmodern}
\usepackage[sfdefault,lining]{FiraSans}
\usepackage[a4paper,top=2.3cm,bottom=2.2cm,left=1.8cm,right=1.8cm,headheight=1.3cm,headsep=0.4cm,footskip=1cm]{geometry}
\usepackage{amsmath,amssymb}
\usepackage{icomma}
\usepackage{graphicx}
\usepackage[table]{xcolor}
\usepackage{tikz}
\usepackage{fontawesome5}
\usepackage{ccicons}
\usepackage{fancyhdr}
\usepackage{lastpage}
\usepackage{tasks}
\usepackage{needspace}
\usepackage[most]{tcolorbox}
\usepackage{array}
\usepackage[hidelinks]{hyperref}
\AtBeginDocument{\hypersetup{pdftitle={Soluciones de la ficha de la unidad 2: Álgebra. 1º Bach.},pdfsubject={Matemáticas CCSS 1º Bach.},pdfkeywords={polinomios, Ruffini, teorema del resto, factorización, fracciones algebraicas, ecuaciones de primer y segundo grado, bicuadradas, radicales, racionales, exponenciales, logarítmicas}}}

\definecolor{navy}{HTML}{1B2748}
\definecolor{accent}{HTML}{E8603A}
\definecolor{teal}{HTML}{1F9E90}
\definecolor{soft}{HTML}{EEF1F8}
\definecolor{muted}{HTML}{5D6685}

\setlength{\parindent}{0pt}
\setlength{\parskip}{3pt}
\settasks{label=\alph*), label-format=\bfseries\color{navy}, label-width=1.3em, item-indent=1.8em, column-sep=0.6em, after-item-skip=5pt, before-skip=3pt, after-skip=4pt}

% ---------- capçalera i peu ----------
\newcommand{\llicencia}{\href{https://creativecommons.org/licenses/by-nc-sa/4.0/deed.es}{\ccbyncsa}\ \ Licencia Creative Commons Atribución-NoComercial-CompartirIgual 4.0 Internacional (CC BY-NC-SA 4.0)}
\pagestyle{fancy}
\fancyhf{}
\fancyhead[L]{\raisebox{-0.35\height}{\includegraphics[height=0.95cm]{logo-gva.png}}}
\fancyhead[C]{\small\color{muted}Soluciones de la ficha · Unidad 2: Álgebra · 1\textsuperscript{o} Bach.}
\fancyhead[R]{\raisebox{-0.35\height}{\includegraphics[height=0.95cm]{logo-vermellar.png}}}
\fancyfoot[L]{\scriptsize\color{muted}\llicencia}
\fancyfoot[R]{\small\color{muted}\thepage}
\renewcommand{\headrulewidth}{0.4pt}
\renewcommand{\headrule}{\color{muted!40}\hrule height \headrulewidth}
\fancypagestyle{primera}{\fancyhf{}\renewcommand{\headrulewidth}{0pt}
  \fancyfoot[L]{\scriptsize\color{muted}\llicencia}
  \fancyfoot[R]{\small\color{muted}\thepage}}

% ---------- nivells i exercicis ----------
\newcommand{\estrelles}[1]{{\color{accent}\ifcase#1\or\faStar\faStar[regular]\faStar[regular]\or\faStar\faStar\faStar[regular]\or\faStar\faStar\faStar\fi}}
\newcounter{exer}
\newcommand{\ex}[1]{\par\needspace{4\baselineskip}\medskip\refstepcounter{exer}\noindent{\color{navy}\bfseries\theexer.}\ \ }

\newcommand{\seccio}[2]{\par\needspace{16\baselineskip}\bigskip
  \begin{tcolorbox}[colback=navy,colframe=navy,coltext=white,arc=2mm,boxrule=0pt,left=3mm,right=3mm,top=1.2mm,bottom=1.2mm]
  {\large\bfseries #1\quad #2}
  \end{tcolorbox}\vspace{-2pt}}

\newcommand{\exf}[1]{\par\medskip\refstepcounter{exer}\noindent{\color{navy}\bfseries\theexer.}\ \ }
\providecommand{\lt}{<}\providecommand{\gt}{>}
\newcolumntype{C}[1]{>{\centering\arraybackslash}p{#1}}
% quadre de recordatori (només a la fitxa d'exercicis)
\newcommand{\recorda}[1]{\par\smallskip\begin{tcolorbox}[colback=accent!7,colframe=accent!55,arc=2mm,boxrule=0.5pt,left=2.5mm,right=2.5mm,top=1mm,bottom=1mm]\small{\color{accent}\faLightbulb[regular]}\ \textbf{Recuerda:}\ #1\end{tcolorbox}}
\newcommand{\res}[1]{\boxed{#1}}\newcommand{\rest}[1]{\fbox{#1}}

\def\fmneed{11}% que el títol d'una secció no quede sol al final del full
% =====================================================================
% Capçalera FentMates de les fitxes: un color per a cada curs.
% Els colors són els de la pàgina inicial (LEVELS a assets/js/materials.js);
% el logo ve de _eines/logo/pdf/ (python3 _eines/logo/genera.py).
%
% Ús, al final del preàmbul de la unitat:
%   % =====================================================================
% Capçalera FentMates de les fitxes: un color per a cada curs.
% Els colors són els de la pàgina inicial (LEVELS a assets/js/materials.js);
% el logo ve de _eines/logo/pdf/ (python3 _eines/logo/genera.py).
%
% Ús, al final del preàmbul de la unitat:
%   % =====================================================================
% Capçalera FentMates de les fitxes: un color per a cada curs.
% Els colors són els de la pàgina inicial (LEVELS a assets/js/materials.js);
% el logo ve de _eines/logo/pdf/ (python3 _eines/logo/genera.py).
%
% Ús, al final del preàmbul de la unitat:
%   \input{../comu/fentmates.tex}
%   \fmcurs{1eso}
% i, al començament del document:
%   \capcalera{Matemàtiques}{1r ESO}{Fitxa d'exercicis · Unitat 2}{Divisibilitat}
% Opcional abans de \input: \def\fmneed{14} (espai mínim abans de cada secció).
% =====================================================================
\makeatletter
\def\fm@logodir{../../logo/pdf}
\providecommand{\fmneed}{10}
\def\fm@cursnom{}

% clau del curs, color i color del text damunt del color
\newcommand{\fm@curs}[3]{\@namedef{fm@color@#1}{#2}\@namedef{fm@text@#1}{#3}}
\fm@curs{1eso}{E6B84A}{navy}
\fm@curs{2eso}{3B82F6}{white}
\fm@curs{3eso}{FF7A50}{white}
\fm@curs{4eso-a}{A78BFA}{white}
\fm@curs{4eso-b}{6D28D9}{white}
\fm@curs{1batx-mi}{9B1C1C}{white}
\fm@curs{1batx-ccss}{EF6B6B}{white}
\fm@curs{2batx-mii}{1E6B47}{white}
\fm@curs{2batx-ccss}{2BB3A3}{white}

\newcommand{\fmcurs}[1]{%
  \@ifundefined{fm@color@#1}{\PackageError{fentmates}{Curs desconegut: #1}{}}{}%
  \edef\fm@tmp{\noexpand\definecolor{curs}{HTML}{\@nameuse{fm@color@#1}}%
               \noexpand\colorlet{curstext}{\@nameuse{fm@text@#1}}}\fm@tmp
  \colorlet{cursfosc}{curs!72!black}%
  \def\fm@logo{\fm@logodir/vertical-#1.pdf}%
  \def\fm@marca{\fm@logodir/marca-#1.pdf}}

% ---------- capçalera de la primera pàgina ----------
% #1 matèria (petit), #2 curs, #3 tipus de document i unitat, #4 títol
\newcommand{\capcalera}[4]{%
  \gdef\fm@cursnom{#2}%
  \par\noindent
  \begin{tcolorbox}[colback=curs,colframe=curs,boxrule=0pt,arc=2.2mm,width=3.5cm,height=2.15cm,
      valign=center,left=3mm,right=2mm,top=1mm,bottom=1mm,nobeforeafter,box align=center]
    \color{curstext}{\footnotesize #1}\par\vspace{1pt}{\fontsize{19}{21}\selectfont\bfseries #2\par}%
  \end{tcolorbox}\hspace{5mm}%
  \begin{minipage}[c]{\dimexpr\linewidth-3.5cm-5mm-2.9cm\relax}
    {\color{cursfosc}\small\bfseries\MakeUppercase{#3}}\par\vspace{2pt}
    {\color{navy}\fontsize{20}{23}\selectfont\bfseries #4\par}
  \end{minipage}\hfill
  \begin{minipage}[c]{2.6cm}\raggedleft\includegraphics[width=2.6cm]{\fm@logo}\end{minipage}\par
  \vspace{2pt}}

% ---------- barres de secció del color del curs ----------
\renewcommand{\seccio}[2]{\par\needspace{\fmneed\baselineskip}\bigskip
  \begin{tcolorbox}[colback=curs,colframe=curs,coltext=curstext,arc=2mm,boxrule=0pt,left=3mm,right=3mm,top=1.2mm,bottom=1.2mm]
  {\large\bfseries #1\quad #2}
  \end{tcolorbox}\vspace{-2pt}}

% ---------- estrelles de dificultat del color del curs (a Batxillerat no se'n posen) ----------
\@ifundefined{estrelles}{}{%
  \renewcommand{\estrelles}[1]{{\color{curs}\ifcase#1\or\faStar\faStar[regular]\faStar[regular]\or\faStar\faStar\faStar[regular]\or\faStar\faStar\faStar\fi}}}

% ---------- capçalera i peu de les pàgines ----------
% primera pàgina: logos institucionals al peu
\fancypagestyle{primera}{\fancyhf{}\renewcommand{\headrulewidth}{0pt}
  \fancyfoot[L]{\raisebox{-0.32\height}{\includegraphics[height=0.78cm]{logo-gva.png}}\hspace{5mm}%
    \raisebox{-0.32\height}{\includegraphics[height=0.78cm]{logo-vermellar.png}}}
  \fancyfoot[R]{\scriptsize\color{muted}\href{https://creativecommons.org/licenses/by-nc-sa/4.0/deed.ca}{\ccbyncsa}\ \ CC BY-NC-SA 4.0 · fentmates.org\qquad\small\thepage}}
% la resta: el símbol a l'esquerra i el curs a la dreta, amb el títol al centre com abans
\fancyhead[L]{\raisebox{-0.3\height}{\includegraphics[height=0.62cm]{\fm@marca}}}
\fancyhead[R]{\small\bfseries\color{cursfosc}\fm@cursnom}
\renewcommand{\headrule}{\color{curs!45}\hrule height \headrulewidth}
\makeatother

%   \fmcurs{1eso}
% i, al començament del document:
%   \capcalera{Matemàtiques}{1r ESO}{Fitxa d'exercicis · Unitat 2}{Divisibilitat}
% Opcional abans de \input: \def\fmneed{14} (espai mínim abans de cada secció).
% =====================================================================
\makeatletter
\def\fm@logodir{../../logo/pdf}
\providecommand{\fmneed}{10}
\def\fm@cursnom{}

% clau del curs, color i color del text damunt del color
\newcommand{\fm@curs}[3]{\@namedef{fm@color@#1}{#2}\@namedef{fm@text@#1}{#3}}
\fm@curs{1eso}{E6B84A}{navy}
\fm@curs{2eso}{3B82F6}{white}
\fm@curs{3eso}{FF7A50}{white}
\fm@curs{4eso-a}{A78BFA}{white}
\fm@curs{4eso-b}{6D28D9}{white}
\fm@curs{1batx-mi}{9B1C1C}{white}
\fm@curs{1batx-ccss}{EF6B6B}{white}
\fm@curs{2batx-mii}{1E6B47}{white}
\fm@curs{2batx-ccss}{2BB3A3}{white}

\newcommand{\fmcurs}[1]{%
  \@ifundefined{fm@color@#1}{\PackageError{fentmates}{Curs desconegut: #1}{}}{}%
  \edef\fm@tmp{\noexpand\definecolor{curs}{HTML}{\@nameuse{fm@color@#1}}%
               \noexpand\colorlet{curstext}{\@nameuse{fm@text@#1}}}\fm@tmp
  \colorlet{cursfosc}{curs!72!black}%
  \def\fm@logo{\fm@logodir/vertical-#1.pdf}%
  \def\fm@marca{\fm@logodir/marca-#1.pdf}}

% ---------- capçalera de la primera pàgina ----------
% #1 matèria (petit), #2 curs, #3 tipus de document i unitat, #4 títol
\newcommand{\capcalera}[4]{%
  \gdef\fm@cursnom{#2}%
  \par\noindent
  \begin{tcolorbox}[colback=curs,colframe=curs,boxrule=0pt,arc=2.2mm,width=3.5cm,height=2.15cm,
      valign=center,left=3mm,right=2mm,top=1mm,bottom=1mm,nobeforeafter,box align=center]
    \color{curstext}{\footnotesize #1}\par\vspace{1pt}{\fontsize{19}{21}\selectfont\bfseries #2\par}%
  \end{tcolorbox}\hspace{5mm}%
  \begin{minipage}[c]{\dimexpr\linewidth-3.5cm-5mm-2.9cm\relax}
    {\color{cursfosc}\small\bfseries\MakeUppercase{#3}}\par\vspace{2pt}
    {\color{navy}\fontsize{20}{23}\selectfont\bfseries #4\par}
  \end{minipage}\hfill
  \begin{minipage}[c]{2.6cm}\raggedleft\includegraphics[width=2.6cm]{\fm@logo}\end{minipage}\par
  \vspace{2pt}}

% ---------- barres de secció del color del curs ----------
\renewcommand{\seccio}[2]{\par\needspace{\fmneed\baselineskip}\bigskip
  \begin{tcolorbox}[colback=curs,colframe=curs,coltext=curstext,arc=2mm,boxrule=0pt,left=3mm,right=3mm,top=1.2mm,bottom=1.2mm]
  {\large\bfseries #1\quad #2}
  \end{tcolorbox}\vspace{-2pt}}

% ---------- estrelles de dificultat del color del curs (a Batxillerat no se'n posen) ----------
\@ifundefined{estrelles}{}{%
  \renewcommand{\estrelles}[1]{{\color{curs}\ifcase#1\or\faStar\faStar[regular]\faStar[regular]\or\faStar\faStar\faStar[regular]\or\faStar\faStar\faStar\fi}}}

% ---------- capçalera i peu de les pàgines ----------
% primera pàgina: logos institucionals al peu
\fancypagestyle{primera}{\fancyhf{}\renewcommand{\headrulewidth}{0pt}
  \fancyfoot[L]{\raisebox{-0.32\height}{\includegraphics[height=0.78cm]{logo-gva.png}}\hspace{5mm}%
    \raisebox{-0.32\height}{\includegraphics[height=0.78cm]{logo-vermellar.png}}}
  \fancyfoot[R]{\scriptsize\color{muted}\href{https://creativecommons.org/licenses/by-nc-sa/4.0/deed.ca}{\ccbyncsa}\ \ CC BY-NC-SA 4.0 · fentmates.org\qquad\small\thepage}}
% la resta: el símbol a l'esquerra i el curs a la dreta, amb el títol al centre com abans
\fancyhead[L]{\raisebox{-0.3\height}{\includegraphics[height=0.62cm]{\fm@marca}}}
\fancyhead[R]{\small\bfseries\color{cursfosc}\fm@cursnom}
\renewcommand{\headrule}{\color{curs!45}\hrule height \headrulewidth}
\makeatother

%   \fmcurs{1eso}
% i, al començament del document:
%   \capcalera{Matemàtiques}{1r ESO}{Fitxa d'exercicis · Unitat 2}{Divisibilitat}
% Opcional abans de \input: \def\fmneed{14} (espai mínim abans de cada secció).
% =====================================================================
\makeatletter
\def\fm@logodir{../../logo/pdf}
\providecommand{\fmneed}{10}
\def\fm@cursnom{}

% clau del curs, color i color del text damunt del color
\newcommand{\fm@curs}[3]{\@namedef{fm@color@#1}{#2}\@namedef{fm@text@#1}{#3}}
\fm@curs{1eso}{E6B84A}{navy}
\fm@curs{2eso}{3B82F6}{white}
\fm@curs{3eso}{FF7A50}{white}
\fm@curs{4eso-a}{A78BFA}{white}
\fm@curs{4eso-b}{6D28D9}{white}
\fm@curs{1batx-mi}{9B1C1C}{white}
\fm@curs{1batx-ccss}{EF6B6B}{white}
\fm@curs{2batx-mii}{1E6B47}{white}
\fm@curs{2batx-ccss}{2BB3A3}{white}

\newcommand{\fmcurs}[1]{%
  \@ifundefined{fm@color@#1}{\PackageError{fentmates}{Curs desconegut: #1}{}}{}%
  \edef\fm@tmp{\noexpand\definecolor{curs}{HTML}{\@nameuse{fm@color@#1}}%
               \noexpand\colorlet{curstext}{\@nameuse{fm@text@#1}}}\fm@tmp
  \colorlet{cursfosc}{curs!72!black}%
  \def\fm@logo{\fm@logodir/vertical-#1.pdf}%
  \def\fm@marca{\fm@logodir/marca-#1.pdf}}

% ---------- capçalera de la primera pàgina ----------
% #1 matèria (petit), #2 curs, #3 tipus de document i unitat, #4 títol
\newcommand{\capcalera}[4]{%
  \gdef\fm@cursnom{#2}%
  \par\noindent
  \begin{tcolorbox}[colback=curs,colframe=curs,boxrule=0pt,arc=2.2mm,width=3.5cm,height=2.15cm,
      valign=center,left=3mm,right=2mm,top=1mm,bottom=1mm,nobeforeafter,box align=center]
    \color{curstext}{\footnotesize #1}\par\vspace{1pt}{\fontsize{19}{21}\selectfont\bfseries #2\par}%
  \end{tcolorbox}\hspace{5mm}%
  \begin{minipage}[c]{\dimexpr\linewidth-3.5cm-5mm-2.9cm\relax}
    {\color{cursfosc}\small\bfseries\MakeUppercase{#3}}\par\vspace{2pt}
    {\color{navy}\fontsize{20}{23}\selectfont\bfseries #4\par}
  \end{minipage}\hfill
  \begin{minipage}[c]{2.6cm}\raggedleft\includegraphics[width=2.6cm]{\fm@logo}\end{minipage}\par
  \vspace{2pt}}

% ---------- barres de secció del color del curs ----------
\renewcommand{\seccio}[2]{\par\needspace{\fmneed\baselineskip}\bigskip
  \begin{tcolorbox}[colback=curs,colframe=curs,coltext=curstext,arc=2mm,boxrule=0pt,left=3mm,right=3mm,top=1.2mm,bottom=1.2mm]
  {\large\bfseries #1\quad #2}
  \end{tcolorbox}\vspace{-2pt}}

% ---------- estrelles de dificultat del color del curs (a Batxillerat no se'n posen) ----------
\@ifundefined{estrelles}{}{%
  \renewcommand{\estrelles}[1]{{\color{curs}\ifcase#1\or\faStar\faStar[regular]\faStar[regular]\or\faStar\faStar\faStar[regular]\or\faStar\faStar\faStar\fi}}}

% ---------- capçalera i peu de les pàgines ----------
% primera pàgina: logos institucionals al peu
\fancypagestyle{primera}{\fancyhf{}\renewcommand{\headrulewidth}{0pt}
  \fancyfoot[L]{\raisebox{-0.32\height}{\includegraphics[height=0.78cm]{logo-gva.png}}\hspace{5mm}%
    \raisebox{-0.32\height}{\includegraphics[height=0.78cm]{logo-vermellar.png}}}
  \fancyfoot[R]{\scriptsize\color{muted}\href{https://creativecommons.org/licenses/by-nc-sa/4.0/deed.ca}{\ccbyncsa}\ \ CC BY-NC-SA 4.0 · fentmates.org\qquad\small\thepage}}
% la resta: el símbol a l'esquerra i el curs a la dreta, amb el títol al centre com abans
\fancyhead[L]{\raisebox{-0.3\height}{\includegraphics[height=0.62cm]{\fm@marca}}}
\fancyhead[R]{\small\bfseries\color{cursfosc}\fm@cursnom}
\renewcommand{\headrule}{\color{curs!45}\hrule height \headrulewidth}
\makeatother

\fmcurs{1batx-ccss}

\begin{document}
\thispagestyle{primera}

\capcalera{Matemáticas CCSS}{1\textsuperscript{o} Bach.}{Soluciones de la ficha · Unidad 2}{\raggedright Álgebra}
\vspace{8pt}

\begin{tcolorbox}[colback=soft,colframe=soft,arc=2mm,left=3mm,right=3mm,top=1.5mm,bottom=1.5mm]
\small Intenta resolver cada ejercicio antes de mirar su solución. Los resultados se dan de manera exacta (fracciones y raíces) y, cuando hace falta, con el valor aproximado. Puedes comprobar las divisiones, la factorización, las fracciones y las ecuaciones de primer y segundo grado con las herramientas interactivas de los apuntes de la unidad.
\end{tcolorbox}

% =====================================================
\seccio{1}{Polinomios. Operaciones}

\exf{1} \begin{tasks}(2)
  \task $\res{2 x^{3} + x^{2} - 3 x + 4}$
  \task $\res{2 x^{3} - 2 x^{2} + 3 x + 1}$
  \task $\res{2 x^{5} - 4 x^{4} + x^{3} + 5 x^{2} - 7 x + 3}$
  \task $\res{x^{4} - 4 x^{3} + 6 x^{2} - 4 x + 1}$
\end{tasks}

\ex{1} \begin{tasks}(1)
  \task $\res{- 16 x^{2} - 17 x + 10}$
  \task $\res{15 x^{2} + 3 x - 3}$
  \task $\res{6 x^{4} - 10 x^{3} - 35 x^{2} + 23 x}$
  \task $\res{3 x^{5} - 2 x^{4} - 20 x^{2} - 16 x + 32}$
  \task $\res{x^{3} - \frac{9 x^{2}}{10} + \frac{4 x}{15}}$
\end{tasks}

\ex{2} $\res{16 x^{4} - 24 x^{3} + 25 x^{2} - 12 x + 4}$

\ex{2} \begin{tasks}(1)
  \task $6x^4+(3b-10)x^3+(2a-5b)x^2+abx$: \ $3b-10=-19$, $ab=-6\Rightarrow \res{a=2,\ b=-3}$
  \task $2x^2+9x+18-3a$: \ $\res{b=9},\ \ 18-3a=3\Rightarrow \res{a=5}$
\end{tasks}

\ex{1} \begin{tasks}(1)
  \task $c(x)=\res{x + \frac{2}{3}},\ \ r(x)=\res{- x - \frac{19}{3}}$
  \task $c(x)=\res{3 x^{2} - 3 x - 5},\ \ r(x)=\res{10 x + 14}$
  \task $c(x)=\res{\frac{1}{2}},\ \ r(x)=\res{- 3 x^{2} + \frac{9}{2}}$
  \task $c(x)=\res{3 x^{2} - 2},\ \ r(x)=\res{3 x + 2}$
\end{tasks}


% =====================================================
\seccio{2}{Regla de Ruffini. Teorema del resto}

\exf{1} \begin{tasks}(1)
  \task ${\def\arraystretch{1.5}\begin{array}{r|rrr|r} & 3 & 2 & 1 & -5\\-2 &  & -6 & 8 & -18\\\hline & 3 & -4 & 9 & -23\end{array}}$\\[2pt]$c(x)=\res{3 x^{2} - 4 x + 9},\ \ r=\res{-23}$
  \task ${\def\arraystretch{1.5}\begin{array}{r|rrrr|r} & 3 & 0 & -2 & -1 & 4\\-2 &  & -6 & 12 & -20 & 42\\\hline & 3 & -6 & 10 & -21 & 46\end{array}}$\\[2pt]$c(x)=\res{3 x^{3} - 6 x^{2} + 10 x - 21},\ \ r=\res{46}$
  \task ${\def\arraystretch{1.5}\begin{array}{r|rrr|r} & 1 & -3 & -1 & 6\\- \frac{3}{2} &  & - \frac{3}{2} & \frac{27}{4} & - \frac{69}{8}\\\hline & 1 & - \frac{9}{2} & \frac{23}{4} & - \frac{21}{8}\end{array}}$\\[2pt]$c(x)=\res{\frac{x^{2}}{2} - \frac{9 x}{4} + \frac{23}{8}},\ \ r=\res{- \frac{21}{8}}\ \ \text{(cociente de Ruffini dividido entre }2)$
  \task ${\def\arraystretch{1.5}\begin{array}{r|rrrr|r} & 4 & 0 & -3 & 0 & 5\\5 &  & 20 & 100 & 485 & 2425\\\hline & 4 & 20 & 97 & 485 & 2430\end{array}}$\\[2pt]$c(x)=\res{- 4 x^{3} - 20 x^{2} - 97 x - 485},\ \ r=\res{2430}\ \ \text{(cociente de Ruffini dividido entre }-1)$
  \task ${\def\arraystretch{1.5}\begin{array}{r|rrrrr|r} & 2 & 0 & -12 & -6 & 0 & -117\\3 &  & 6 & 18 & 18 & 36 & 108\\\hline & 2 & 6 & 6 & 12 & 36 & -9\end{array}}$\\[2pt]$c(x)=\res{2 x^{4} + 6 x^{3} + 6 x^{2} + 12 x + 36},\ \ r=\res{-9}$
\end{tasks}

\ex{1} $P(-1)=1+2-3=\res{0}$, \ $P(3)=9-6-3=\res{0}$. \rest{No}: un polinomio de grado 2 tiene como máximo dos raíces, y ya las hemos encontrado: $P(x)=(x+1)(x-3)$.

\ex{1} \begin{tasks}(5)
  \task $P(1)=0$: \rest{sí}
  \task $P(2)=36$: \rest{no}
  \task $P(-1)=-18$: \rest{no}
  \task $P(-4)=0$: \rest{sí}
  \task $P(-3)=-44$: \rest{no}
\end{tasks}

\ex{1} $M(2)=16-20+8-4=0$: \rest{sí es divisible}.\par Por Ruffini: ${\def\arraystretch{1.5}\begin{array}{r|rrr|r} & 2 & -5 & 4 & -4\\2 &  & 4 & -2 & 4\\\hline & 2 & -1 & 2 & 0\end{array}}\Rightarrow \res{N(x)=2x^2-x+2}$

\ex{1} \begin{tasks}(2)
  \task $P(-2)=-8+4+5=\res{1}$
  \task $P(-1)=1+2+1=\res{4}$
  \task $P(1)=2-3+1-1=\res{-1}$
  \task $P(2)=32-32=\res{0}$ (exacta)
\end{tasks}

\ex{2} \begin{tasks}(1)
  \task $P(2)=k + 8=0\Rightarrow \res{k=-8}$
  \task $P(-2)=- 2 k - 12=0\Rightarrow \res{k=-6}$
  \task $P(-4)=- 63 k + 252=0\Rightarrow \res{k=4}$
  \task $P(3)=- 6 k + 36=0\Rightarrow \res{k=6}$
\end{tasks}

\ex{2} \begin{tasks}(1)
  \task $P(-3)=a - 3=0\Rightarrow \res{a=3}$
  \task $P(- \frac{3}{2})=\frac{3 a}{2} - \frac{3}{2}=0\Rightarrow \res{a=1}$
  \task $P(\frac{4}{3})=\frac{16 a}{9} + \frac{160}{9}=0\Rightarrow \res{a=-10}$
  \task $P(5)=b + 2930=0\Rightarrow \res{b=-2930}$
\end{tasks}($P(r)$ tiene que ser 0, donde $r$ es la raíz del divisor; es lo mismo que obtener resto 0 en Ruffini.)

\ex{2} $P(-3)=-27-9(3k+2)-3=-48-27k$.\begin{tasks}(2)
  \task $-48-27k=0\Rightarrow \res{k=-\frac{16}{9}}$
  \task $-48-27k=15\Rightarrow \res{k=-\frac73}$
\end{tasks}

\ex{2} Teorema del resto: $P(-2)=32-72+8+12+3m=-20+3m=12\Rightarrow \res{m=\frac{32}{3}}$


% =====================================================
\seccio{3}{Factorización de polinomios}

\exf{1} \begin{tasks}(1)
  \task $\res{x=-5,\ x=2,\ x=3}$
  \task $\res{x=- \frac{1}{2},\ x=0,\ x=2\ (\text{triple})}$
  \task $\res{x=-1\ (\text{triple}),\ x=- \frac{2}{3},\ x=\frac{1}{2}}$
  \task $\res{x=-2,\ x=1,\ x=4}$ \ $\big((x - 4) (x - 1) (x + 2)\big)$
  \task $\res{x=-5,\ x=-2,\ x=-1}$ \ $\big((x + 1) (x + 2) (x + 5)\big)$
  \task $\res{x=0,\ x=\frac{1}{2},\ x=1,\ x=2\ (\text{doble})}$ \ $\big(x (x - 2)^{2} (x - 1) (2 x - 1)\big)$
\end{tasks}

\ex{2} \begin{tasks}(1)
  \task $\res{(x - 3) (x - 1) (x + 2)}$ \ \small(raíces: $x=-2,\ x=1,\ x=3$)
  \task $\res{(x - 1)^{2} (x + 3)}$ \ \small(raíces: $x=-3,\ x=1\ (\text{doble})$)
  \task $\res{(x - 1) (x + 1)^{2} (2 x - 3)}$ \ \small(raíces: $x=-1\ (\text{doble}),\ x=1,\ x=\frac{3}{2}$)
  \task $\res{(x - 2) (2 x - 1) (2 x + 3)}$ \ \small(raíces: $x=- \frac{3}{2},\ x=\frac{1}{2},\ x=2$)
  \task $\res{x^{2} (x - 2) (x + 2)}$ \ \small(raíces: $x=-2,\ x=0\ (\text{doble}),\ x=2$)
  \task $\res{x (x + 1)^{2} (2 x - 1) (3 x + 1)}$ \ \small(raíces: $x=-1\ (\text{doble}),\ x=- \frac{1}{3},\ x=0,\ x=\frac{1}{2}$)
  \task $\res{(x - 2) (x - 1) (x + 1) (x + 2) (x^{2} + 1)}$ \ \small(raíces: $x=-2,\ x=-1,\ x=1,\ x=2$)
  \task $\res{3 (x - 2)^{2}}$ \ \small(raíces: $x=2\ (\text{doble})$)
  \task $t=x^2$: $\frac{t^2}{4}+\frac34t+\frac9{16}=0\Rightarrow\Delta=\frac9{16}-\frac9{16}=0$, $t=-\frac32\lt0$: \rest{no tiene raíces reales}. $P(x)=\res{\frac{1}{16}(2x^2+3)^2}$
\end{tasks}

\ex{2} $P(x)=ax^2+bx+c$: \ $c=-3$, \ $a+b-3=-6$, \ $9a+3b-3=0\Rightarrow a=2,\ b=-5$: \ $\res{P(x)=2x^2-5x-3=(x-3)(2x+1)}$

\ex{1} $P(2)=8m-24-8+8=8m-24=0\Rightarrow \res{m=3}$; \ $3x^3-6x^2-4x+8=\res{(x - 2) (3 x^{2} - 4)}$

\ex{2} $P(2)=0$: $4a+2b=-2$; \ $P(-3)=0$: $9a-3b=33$. \ $\res{a=2,\ b=-5}$ \ (y entonces $P(x)=(x-2)(x+3)(x+1)$).

\ex{2} $P(x)=a(x-1)(x+2)(x-3)$; \ $P(0)=6a=12\Rightarrow a=2$: \ $\res{2 x^{3} - 4 x^{2} - 10 x + 12}$


% =====================================================
\seccio{4}{Fracciones algebraicas}

\exf{2} \begin{tasks}(1)
  \task $\dfrac{2 (x - 2) (x + 2)}{x (x - 1) (x + 2)}=\res{\frac{2 (x - 2)}{x (x - 1)}}$
  \task $\dfrac{x (x - 1) (x + 1)}{x (x + 1) (x + 2)}=\res{\frac{x - 1}{x + 2}}$
  \task $\dfrac{x^{2} (x - 3) (x + 1)}{x^{2} (x - 3) (x + 3)}=\res{\frac{x + 1}{x + 3}}$
  \task $\dfrac{(x - 1)^{3}}{x (x - 1)^{2}}=\res{\frac{x - 1}{x}}$
  \task $\dfrac{(x - 1)^{2} (x + 1) (2 x - 3)}{(x - 1) (x + 2) (2 x - 3)}=\res{\frac{(x - 1) (x + 1)}{x + 2}}$
  \task $\dfrac{(x - 3) (x^{2} + 3 x + 9)}{(x - 3) (x + 3)}=\res{\frac{x^{2} + 3 x + 9}{x + 3}}$
  \task $\dfrac{(x - 3) (x - 2) (x + 3)}{(x - 3) (x - 2)^{2}}=\res{\frac{x + 3}{x - 2}}$
  \task $\dfrac{(x - 2) (x - 1)^{3}}{(x - 2)^{2} (x - 1)^{2}}=\res{\frac{x - 1}{x - 2}}$
\end{tasks}

\ex{2} \begin{tasks}(2)
  \task $\res{\frac{x + 2}{(x - 1) (x + 1)}}$
  \task $\res{\frac{3 x + 1}{x (x + 1)}}$
  \task $\res{\frac{2 x - 5}{4 x (x - 1)}}$
  \task $\res{\frac{2 (2 x + 1)}{x (x - 1) (x + 1)}}$
  \task $\res{\frac{5 x^{2} - 3 x + 1}{(x - 2) (x + 2)}}$
  \task $\res{\frac{2 (x^{2} + x - 1)}{(x - 1)^{2} (x + 1)}}$
\end{tasks}

\ex{2} \begin{tasks}(2)
  \task $\res{\frac{1}{4}}$
  \task $\res{\frac{x + 3}{(x - 2) (x + 1)}}$
  \task $\res{\frac{(3 x + 2) (4 x + 1)}{(x - 2) (x + 2)}}$
  \task $\res{- \frac{2}{3 x^{2}}}$
  \task $\res{\frac{(x - 2) (x + 5)}{x + 1}}$
  \task $\res{\frac{x + 2}{x}}$
\end{tasks}

\ex{3} \begin{tasks}(2)
  \task $\res{- \frac{2 x^{2} - 3 x - 3}{x - 1}}$
  \task $\res{\frac{x^{2} - 6 x - 1}{2 (x + 1)^{2}}}$
  \task $\res{x}$
  \task $\res{1}$
  \task $\res{\frac{3 x (x + 1)}{x^{2} + 1}}$
  \task $\res{\frac{2 (5 x + 6)}{(x - 2) (x + 2)}}$
  \task $\res{\frac{x}{x + 1}}$
  \task $\res{\frac{3 x + 2}{2 x + 1}}$
\end{tasks}


% =====================================================
\seccio{5}{Ecuaciones polinómicas}

\exf{1} \begin{tasks}(2)
  \task $5 x=18$: \ $\res{x=\frac{18}{5}}$
  \task $3 x=0$: \ $\res{x=0}$
  \task $4 x=8$: \ $\res{x=2}$
  \task $0\cdot x=2$: \ $\res{\text{no tiene solución}}$
  \task $0\cdot x=0$: \rest{identidad} (infinitas soluciones)
  \task $- 2 x=2$: \ $\res{x=-1}$
\end{tasks}

\ex{2} \begin{tasks}(1)
  \task m.c.m.$\,=20$: \ $- 32 x=-9$: \ $\res{x=\frac{9}{32}}$
  \task m.c.m.$\,=12$: \ $- 5 x=-23$: \ $\res{x=\frac{23}{5}}$
  \task m.c.m.$\,=6$: \ $x=13$: \ $\res{x=13}$
  \task m.c.m.$\,=6$: \ $0\cdot x=5$: \ $\res{\text{no tiene solución}}$
  \task m.c.m.$\,=12$: \ $4 x=-5$: \ $\res{x=- \frac{5}{4}}$
  \task m.c.m.$\,=10$: \ $2 x=11$: \ $\res{x=\frac{11}{2}}$
\end{tasks}

\ex{1} \begin{tasks}(2)
  \task $\res{x=2,\ x=3}$
  \task $\res{x=-2,\ x=\frac{1}{2}}$
  \task $\res{x=-2,\ x=2}$
  \task $\Delta\lt0$: $\res{\text{no tiene solución real}}$
  \task $\res{x=\frac{1}{2}\ (\text{doble})}$
  \task $\res{x=0,\ x=2}$
  \task $x^{2} + 5 x=0$: \ $\res{x=-5,\ x=0}$
  \task Por 12: \ $6x^2+6-6x+9+2x^2=59\Rightarrow8x^2-6x-44=0\Rightarrow4x^2-3x-22=0$: \ $\res{x=-2,\ x=\frac{11}{4}}$
  \task $6x^2+2=7\Rightarrow x^2=\frac56$: \ $\res{x=- \frac{\sqrt{30}}{6},\ x=\frac{\sqrt{30}}{6}}$
  \task $3 x^{2} - 5 x - 2=0$: \ $\res{x=- \frac{1}{3},\ x=2}$
\end{tasks}

\ex{2} \begin{tasks}(1)
  \task $t=x^2$: $t^2-17t+16=0\Rightarrow t=16,\ t=1$: \ $\res{x=\pm4,\ x=\pm1}$
  \task $t=x^3$: $t^2-9t+8=0\Rightarrow t=8,\ t=1$: \ $\res{x=2,\ x=1}$
  \task $t^2-8t-9=0\Rightarrow t=9$ ($t=-1$ no da soluciones): \ $\res{x=\pm3}$
  \task Multiplicamos por $x^2$: $x^4-10x^2+9=0\Rightarrow t=9,\ t=1$: \ $\res{x=\pm3,\ x=\pm1}$
  \task Desarrollando: $2x^4+12x^2-14=0\Rightarrow t^2+6t-7=0\Rightarrow t=1$ ($t=-7$ no da soluciones): \ $\res{x=\pm1}$
  \task $4t^2-17t+4=0\Rightarrow t=4,\ t=\frac14$: \ $\res{x=\pm2,\ x=\pm\frac12}$
\end{tasks}

\ex{2} \begin{tasks}(1)
  \task $(x - 1) (x + 2) (2 x + 3) (3 x - 1)=0$: \ $\res{x=-2,\ x=- \frac{3}{2},\ x=\frac{1}{3},\ x=1}$
  \task $(x - 2) (2 x + 3) (x^{2} + 3)=0$: \ $\res{x=- \frac{3}{2},\ x=2}$
  \task $x (x + 1) (2 x - 1) (3 x + 1)=0$: \ $\res{x=-1,\ x=- \frac{1}{3},\ x=0,\ x=\frac{1}{2}}$
  \task $(x - 2) (x - 1) (x + 3)=0$: \ $\res{x=-3,\ x=1,\ x=2}$
  \task $(x - 3)^{2} (x + 3)=0$: \ $\res{x=-3,\ x=3}$
  \task $(x - 3) (x + 2) (x^{2} + 1)=0$: \ $\res{x=-2,\ x=3}$
  \task $(x - 1) (x + 1) (x^{2} + 1)=0$: \ $\res{x=-1,\ x=1}$
  \task $x (x - 3) (x + 1) (x + 3)=0$: \ $\res{x=-3,\ x=-1,\ x=0,\ x=3}$
\end{tasks}

\ex{1} \begin{tasks}(2)
  \task $\res{x=-2,\ x=- \sqrt{3},\ x=\sqrt{3},\ x=2}$
  \task $\res{x=-3,\ x=- \frac{2}{3},\ x=0,\ x=3,\ x=4}$
  \task $\res{x=-3,\ x=2}$
  \task $\res{x=-1,\ x=2,\ x=\frac{7}{2}}$
\end{tasks}


% =====================================================
\seccio{6}{Ecuaciones con radicales}

\exf{2} \begin{tasks}(1)
  \task $x-6=\sqrt x\Rightarrow x^2-13x+36=0\Rightarrow x=9,\ x=4$; \ $x=4$ no vale ($4-2\neq6$): \ $\res{x=9}$
  \task $2x-3=(1-x)^2\Rightarrow(x-2)^2=0$; \ $x=2$: $-1+1=0\neq2$: \ $\res{\text{no tiene solución}}$
  \task $2x^2+2x-3=(x+1)^2\Rightarrow x^2=4$; \ $x=2$ no vale ($-1-2\lt0$): \ $\res{x=-2}$
  \task $9(7x+1)=(5x-7)^2\Rightarrow25x^2-133x+40=0\Rightarrow x=5,\ x=\frac{8}{25}$; \ $x=\frac8{25}$ no vale: \ $\res{x=5}$
  \task $2x-1=12$: \ $\res{x=\frac{13}{2}}$
  \task $2-5x=3x^2\Rightarrow3x^2+5x-2=0\Rightarrow x=-2,\ x=\frac13$; \ $x=\frac13$ no vale: \ $\res{x=-2}$
  \task $2+\sqrt{x-4}=12-x\Rightarrow x-4=(10-x)^2\Rightarrow x^2-21x+104=0\Rightarrow x=8,\ x=13$; \ $x=13$ no vale: \ $\res{x=8}$
  \task $x=4-4\sqrt{x-4}+x-4\Rightarrow\sqrt{x-4}=0$: \ $\res{x=4}$
  \task $x-46=-14\sqrt{x+5}\Rightarrow x^2-288x+1136=0\Rightarrow x=4,\ x=284$; \ $x=284$ no vale: \ $\res{x=4}$
  \task $10\sqrt{x-1}=20\Rightarrow\sqrt{x-1}=2$: \ $\res{x=5}$
  \task $x+3=1+2\sqrt{x-2}+x-2\Rightarrow\sqrt{x-2}=2$: \ $\res{x=6}$ \ ($3-2=1$ \checkmark)
  \task $\sqrt{1-4x}=3\sqrt{3+x}\Rightarrow1-4x=27+9x$: \ $\res{x=-2}$
\end{tasks}

\ex{3} \begin{tasks}(1)
  \task $x^2-1=(x-1)^3\Rightarrow x(x-1)(x-3)=0$; todas valen ($\sqrt[3]{-1}=-1$): \ $\res{x=0,\ x=1,\ x=3}$
  \task $4x-1=(x-4)^3\Rightarrow(x-7)(x^2-5x+9)=0$: \ $\res{x=7}$
  \task $(2x+2)^2=6x+10\Rightarrow2x^2+x-3=0\Rightarrow x=1,\ x=-\frac32$; \ $x=-\frac32$ no vale: \ $\res{x=1}$
\end{tasks}

\ex{2} $\sqrt{x+1}=2\sqrt{x-11}\Rightarrow x+1=4x-44$: \ $\res{x=15}$ \ (comprobación: $\sqrt{16}=4=2\sqrt4$).


% =====================================================
\seccio{7}{Ecuaciones racionales}

\exf{2} \begin{tasks}(1)
  \task $8x^2+8x+8=7x^3\Rightarrow(x-2)(7x^2+6x+4)=0$ ($\Delta\lt0$): \ $\res{x=2}$
  \task Multiplicamos por $3(x-1)$: $2x^2+4x+9=11\Rightarrow x^2+2x-1=0$: \ $\res{x=-1\pm\sqrt2}$
  \task $2x(x+2)+3x(x-2)=6x^2\Rightarrow x^2+2x=0\Rightarrow x=0,\ x=-2$; \ $x=-2$ anula un denominador: \ $\res{x=0}$
  \task m.c.m. $=(x-1)(x-2)(x-4)$: $3(x-1)+(x-5)(x-2)=(x-1)(x-2)\Rightarrow x^2-4x+7=x^2-3x+2$: \ $\res{x=5}$
  \task Multiplicamos por $\sqrt2\,x$: $x^2+2=2x^2\Rightarrow$ $\res{x=\pm\sqrt2}$
  \task $-\dfrac{x}{x-2}=-x\Rightarrow x=x(x-2)\Rightarrow x^2-3x=0$: \ $\res{x=0,\ x=3}$
  \task $2x^2+2(x+1)^2=5x(x+1)\Rightarrow x^2+x-2=0$: \ $\res{x=1,\ x=-2}$ (ninguna anula los denominadores)
  \task Simplificando: $\dfrac{2x(x+1)}{(x-1)(x+3)}=2\Rightarrow x^2+x=x^2+2x-3$: \ $\res{x=3}$
  \task El primer miembro es $\dfrac{x+1}{x}$: es una \rest{identidad}, se cumple para todo $x\neq0,\ x\neq-1$.
\end{tasks}

\ex{2} m.c.m. $=(x+2)(x-2)$: \ $(x-2)^2-(x+2)=2x-6+(x^2-4)\Rightarrow x^2-5x+2=x^2+2x-10\Rightarrow7x=12$, y al comprobarlo: $x=\frac{12}{7}$ sí es válida: \ $\res{x=\frac{12}{7}}$

\ex{3} $\dfrac{600}{x-10}=\dfrac{600}{x}+5\Rightarrow x^2-10x-1200=0\Rightarrow x=40$ ($x=-30$ no tiene sentido): \ $\res{40\ \text{personas}}$


% =====================================================
\seccio{8}{Ecuaciones exponenciales}

\exf{2} \begin{tasks}(1)
  \task $2^{2x^2+2}=2^{5x+5}\Rightarrow2x^2-5x-3=0$: \ $\res{x=3,\ x=-\frac12}$
  \task $262144=4^9\Rightarrow(x-2)^2=9$: \ $\res{x=5,\ x=-1}$
  \task $x+4=\dfrac{\log5}{\log3}$: \ $\res{x=\frac{\log5}{\log3}-4\approx-2,535}$
  \task $2^{-x}=2^{2x(x-1)}\Rightarrow2x^2-x=0$: \ $\res{x=0,\ x=\frac12}$
  \task $13^{x^2+2x}=13^{-1}\Rightarrow(x+1)^2=0$: \ $\res{x=-1}$
  \task $2x^2-4x-2=4\Rightarrow x^2-2x-3=0$: \ $\res{x=3,\ x=-1}$
  \task $6^{4x-3}=12\Rightarrow4x-3=\dfrac{\log12}{\log6}$: \ $\res{x=\frac{3+\log_612}{4}\approx1,097}$
  \task $\left(\dfrac32\right)^x=\dfrac32$: \ $\res{x=1}$
  \task $3^{3/x}=3^{x+2}\Rightarrow x^2+2x-3=0$: \ $\res{x=1,\ x=-3}$
  \task $2^{x+4}=2^{3x}\Rightarrow x+4=3x$: \ $\res{x=2}$
  \task $x^2-5x+6=0$: \ $\res{x=2,\ x=3}$
  \task $7^{\frac{2(x-5)}{3}}=7^{-\frac{x-1}{2}}\Rightarrow4(x-5)=-3(x-1)$: \ $\res{x=\frac{23}{7}}$
\end{tasks}

\ex{2} \begin{tasks}(1)
  \task $2^x\left(2+1+\frac12\right)=28\Rightarrow2^x=8$: \ $\res{x=3}$
  \task $5^x\left(125-\frac15\right)=3120\Rightarrow5^x=25$: \ $\res{x=2}$
  \task $4^x\cdot\frac{21}{16}=336\Rightarrow4^x=256$: \ $\res{x=4}$
  \task $3^x\left(9+3+1+\frac13\right)=120\Rightarrow3^x=9$: \ $\res{x=2}$
  \task $2^x\cdot\frac{15}{16}=960\Rightarrow2^x=1024$: \ $\res{x=10}$
  \task $t=6^x$: $t^2+8t-9=0\Rightarrow t=1$ ($t=-9$ no vale): \ $\res{x=0}$
\end{tasks}

\ex{3} \begin{tasks}(1)
  \task $t=3^x$: $t^2+5t-24=0\Rightarrow t=3$ ($t=-8$ no vale): \ $\res{x=1}$
  \task $t=2^x$: $2t^2-5t+2=0\Rightarrow t=2,\ t=\frac12$: \ $\res{x=1,\ x=-1}$
  \task $t=3^x$: $3t^2-10t+3=0\Rightarrow t=3,\ t=\frac13$: \ $\res{x=1,\ x=-1}$
  \task $t=10^{x+2}$: $2t^2+3t-5=0\Rightarrow t=1$ ($t=-\frac52$ no vale) $\Rightarrow x+2=0$: \ $\res{x=-2}$
  \task $t=3^x$: $9t+\frac{t^2}{9}=90\Rightarrow t^2+81t-810=0\Rightarrow t=9$ ($t=-90$ no vale): \ $\res{x=2}$
  \task $t=2^x$: $\frac{t^2}{16}-\frac{5t}{8}+1=0\Rightarrow t^2-10t+16=0\Rightarrow t=8,\ t=2$: \ $\res{x=3,\ x=1}$
  \task $t=2^x$: $4t^2+8t-320=0\Rightarrow t^2+2t-80=0\Rightarrow t=8$ ($t=-10$ no vale): \ $\res{x=3}$
\end{tasks}

\ex{2} $5000\cdot1,04^t=8000\Rightarrow1,04^t=1,6\Rightarrow t=\dfrac{\log1,6}{\log1,04}\approx11,98$: \ $\res{12\ \text{años}}$

\ex{2} \begin{tasks}(1)
  \task $C\cdot1,03^t=2C\Rightarrow t=\dfrac{\log2}{\log1,03}\approx23,45$: \ $\res{24\ \text{años}}$
  \task $1,025^t=1,2\Rightarrow t=\dfrac{\log1,2}{\log1,025}\approx7,38$: \ $\res{8\ \text{años}}$
\end{tasks}


% =====================================================
\seccio{9}{Ecuaciones logarítmicas}

\exf{2} \begin{tasks}(1)
  \task $2(x+3)=x+5$: \ $\res{x=-1}$
  \task $15-2x=x^2\Rightarrow x=3,\ x=-5$; \ $x=-5$ no vale: \ $\res{x=3}$
  \task $\dfrac{2x+3}{x-1}=36\Rightarrow34x=39$: \ $\res{x=\frac{39}{34}}$
  \task $x^2-3x+2=x^2\Rightarrow$ $\res{x=\frac23}$ (vale: $x\gt0$ y $x^2-3x+2=\frac49\gt0$)
  \task $\dfrac{2x-2}{x}=\dfrac{(x-1)^2}{x}\Rightarrow(x-1)(x-3)=0$; \ $x=1$ no vale: \ $\res{x=3}$
  \task $\log x^2=\log\dfrac{x}{20}\Rightarrow x^2=\dfrac x{20}$; \ $x=0$ no vale: \ $\res{x=\frac1{20}}$
  \task $2x(4-5x)=3(2x-x^2)\Rightarrow7x^2-2x=0$; \ $x=0$ no vale: \ $\res{x=\frac27}$
  \task $x^2-1=8\Rightarrow x=\pm3$; \ $x=-3$ no vale: \ $\res{x=3}$
  \task $\dfrac{5(x+1)(x+2)}{x+2}=4\Rightarrow5x+5=4$: \ $\res{x=-\frac15}$
  \task $4-x=(x+2)^2\Rightarrow x^2+5x=0$; \ $x=-5$ no vale: \ $\res{x=0}$
\end{tasks}

\ex{1} $\log x+\log2=1\Rightarrow\log(2x)=\log10$: \ $\res{x=5}$

\ex{1} $2\log x=\log4+\log9\Rightarrow x^2=36$ ($x\gt0$): \ $\res{x=6}$

\ex{2} $\dfrac{x^2}{x+6}=3\Rightarrow x^2-3x-18=0\Rightarrow x=6,\ x=-3$; \ $x=-3$ no vale: \ $\res{x=6}$


% =====================================================
\seccio{10}{Problemas}

\exf{1} $38+x=2(10+x)\Rightarrow$ $\res{x=18\ \text{años}}$

\ex{1} $x+(x+1)+(x+2)=93\Rightarrow3x=90$: \ $\res{30,\ 31,\ 32}$

\ex{1} $1,21x=60,50\Rightarrow$ $\res{x=50\ \text{€}}$

\ex{2} $0,8\cdot0,9\,x=0,72x=36\Rightarrow$ $\res{x=50\ \text{€}}$; \ rebaja total: $\res{28\,\%}$ (no el 30\,\%).

\ex{2} $5x+2\cdot50+0\cdot50=4(x+100)\Rightarrow$ $\res{x=300\ \text{L}}$

\ex{1} $30\cdot5+10x=7(30+x)\Rightarrow3x=60$: \ $\res{20\ \text{L}}$

\ex{2} En 1 kg, $x$ kg de té blanco: $18x+14(1-x)=16,4\Rightarrow x=0,6$: \ $\res{60\,\%\ \text{blanco y }40\,\%\ \text{verde}}$

\ex{1} El 10\,\% son 15: \ $\res{150\ \text{alumnos}}$

\ex{1} África es el 15\,\%: $0,15x=3\Rightarrow$ $\res{20\ \text{millones de euros}}$

\ex{2} $5x-2(20-x)=58\Rightarrow7x=98$: \ $\res{14\ \text{aciertos y }6\ \text{fallos}}$

\ex{3} $(x-20)\left(\dfrac{125}{x}+0,4\right)=147\Rightarrow x^2-75x-6250=0\Rightarrow x=125$ ($x=-50$ no tiene sentido): \ $\res{125\ \text{kg}}$

\ex{2} $x(x+10)=1200\Rightarrow x^2+10x-1200=0\Rightarrow x=30$ ($x=-40$ no tiene sentido): \ $\res{30\ \text{m}\times40\ \text{m}}$

\ex{2} $x^2+(70-x)^2=2500\Rightarrow x^2-70x+1200=0$: \ $\res{40\ \text{m}\times30\ \text{m}}$

\ex{1} $\frac15+\frac1{20}=\frac14$: \ $\res{T=4\ \text{h}}$

\ex{2} $\dfrac1t+\dfrac1{2t}=\dfrac1{10}\Rightarrow t=15$: \ $\res{15\ \text{h y }30\ \text{h}}$

\ex{2} \begin{tasks}(1)
  \task $\res{y(0)=300\ \text{enfermos}}$
  \task $200e^{0,2t}=800\Rightarrow t=5\ln4\approx6,93$: \ $\res{\text{el 7 de enero (hacia el final del día)}}$
\end{tasks}

\ex{2} $2^t=\dfrac{100}{9}\Rightarrow t=\dfrac{\log(100/9)}{\log2}\approx\res{3,47\ \text{h}\approx3\ \text{h}\ 28\ \text{min}}$


\end{document}
